% GENERATED FILE. Edit canonical lessons, then run npm run generate:books.
% canonical-source-hash: fnv1a64:47d4c5549f48d20c
% canonical-lessons: ES-C532-misterio, ES-C532-enigma, ES-C532-ocultar, ES-C532-esconder, ES-C532-secreto

\chapter{Un misterio}
\label{ch:es-un-misterio}
\hlchaptermodality{532}

\begin{chapteropening}
\textbf{By the end of this chapter:} \emph{I can say that there was a mystery or a secret, and that someone hid something.}

Complete the last lesson of chapter 532: I can say that there was a mystery or a secret, and that someone hid something.
\end{chapteropening}

\section[el misterio]{el misterio — the mystery}

\label{lesson:ES-C532-misterio}



\begin{quote}
Before the new one: say the Spanish for to find, then the Spanish for the find.
\end{quote}

\subsection*{You'll want to know: el misterio}
\textbf{el misterio} — \textquotedblleft{}the mystery\textquotedblright{}.

\begin{quote}
Lo que pasó aquella noche sigue siendo un misterio.
\end{quote}

\begin{quote}
\emph{What happened that night is still a mystery.}
\end{quote}

\begin{cousinweb}
\emph{misterio} comes from Latin \emph{mystērium}, a learned word taken from Greek \emph{myst\'{\={e}}rion} 'a secret rite' (EtymDB).
\end{cousinweb}

\begin{grammarlens}[title={What you've built}]
One more word for telling, explaining and describing at B1.
\end{grammarlens}

\subsection*{Your turn}
\begin{itemize}
\raggedright
  \item \emph{Say it:} \emph{el misterio}
  \item \emph{Say it:} \emph{el misterio}, once more
  \item \emph{Say it:} a sentence of your own with \emph{el misterio}
  \item \emph{Recall:} say the Spanish for to find, then the Spanish for the find, then say \emph{el misterio} again
\end{itemize}

\begin{tcolorbox}[breakable,colback=teal!4,colframe=teal!35!black,title={Before you move on}]
What does el misterio mean? (\textquotedblleft{}The mystery\textquotedblright{}.) Say it once more.
\end{tcolorbox}
\section[el enigma]{el enigma — the enigma, the puzzle}

\label{lesson:ES-C532-enigma}



\begin{quote}
Before the new one: say the Spanish for the find, then the Spanish for the mystery.
\end{quote}

\subsection*{You'll want to know: el enigma}
\textbf{el enigma} — \textquotedblleft{}the enigma, the puzzle\textquotedblright{}.

\begin{quote}
Su desaparición sigue siendo un enigma.
\end{quote}

\begin{quote}
\emph{His disappearance is still a puzzle.}
\end{quote}

Masculine although it ends in \emph{-a}: \emph{un enigma}.

\begin{cousinweb}
\emph{enigma} is a learned borrowing from Latin \emph{aenigma}, which comes from Greek \emph{aínigma} 'a riddle' (EtymDB).
\end{cousinweb}

\begin{grammarlens}[title={What you've built}]
One more word for telling, explaining and describing at B1.
\end{grammarlens}

\subsection*{Your turn}
\begin{itemize}
\raggedright
  \item \emph{Say it:} \emph{el enigma}
  \item \emph{Say it:} \emph{el enigma}, once more
  \item \emph{Say it:} a sentence of your own with \emph{el enigma}
  \item \emph{Recall:} say the Spanish for the find, then the Spanish for the mystery, then say \emph{el enigma} again
\end{itemize}

\begin{tcolorbox}[breakable,colback=teal!4,colframe=teal!35!black,title={Before you move on}]
What does el enigma mean? (\textquotedblleft{}The enigma, the puzzle\textquotedblright{}.) Say it once more.
\end{tcolorbox}
\section[ocultar]{ocultar — to conceal, to hide}

\label{lesson:ES-C532-ocultar}



\begin{quote}
Before the new one: say the Spanish for the mystery, then the Spanish for the enigma.
\end{quote}

\subsection*{You'll want to know: ocultar}
\textbf{ocultar} — \textquotedblleft{}to conceal, to hide\textquotedblright{}.

\begin{quote}
Ocultó la verdad durante años.
\end{quote}

\begin{quote}
\emph{She hid the truth for years.}
\end{quote}

A little formal; often for information or feelings: \emph{ocultar la verdad}.

\begin{cousinweb}
\emph{ocultar} is a learned borrowing from Latin \emph{occultāre} 'to hide' (EtymDB), built on \emph{occultus} 'hidden', from \emph{occulere} 'to cover over' (Lewis \& Short).
\end{cousinweb}

\begin{grammarlens}[title={What you've built}]
One more word for telling, explaining and describing at B1.
\end{grammarlens}

\subsection*{Your turn}
\begin{itemize}
\raggedright
  \item \emph{Say it:} \emph{ocultar}
  \item \emph{Say it:} \emph{ocultar}, once more
  \item \emph{Say it:} a sentence of your own with \emph{ocultar}
  \item \emph{Recall:} say the Spanish for the mystery, then the Spanish for the enigma, then say \emph{ocultar} again
\end{itemize}

\begin{tcolorbox}[breakable,colback=teal!4,colframe=teal!35!black,title={Before you move on}]
What does ocultar mean? (\textquotedblleft{}To conceal, to hide\textquotedblright{}.) Say it once more.
\end{tcolorbox}
\section[esconder]{esconder — to hide}

\label{lesson:ES-C532-esconder}



\begin{quote}
Before the new one: say the Spanish for the enigma, then the Spanish for to conceal.
\end{quote}

\subsection*{You'll want to know: esconder}
\textbf{esconder} — \textquotedblleft{}to hide\textquotedblright{}.

\begin{quote}
Escondió el dinero debajo de la cama.
\end{quote}

\begin{quote}
\emph{He hid the money under the bed.}
\end{quote}

\emph{Esconderse} is to hide yourself.

\begin{cousinweb}
\emph{esconder} is inherited from Latin \emph{abscondere} 'to put away, to hide' (Lewis \& Short), built on \emph{abs} 'away' and \emph{condere} 'to put, to store'.
\end{cousinweb}

\begin{grammarlens}[title={What you've built}]
One more word for telling, explaining and describing at B1.
\end{grammarlens}

\subsection*{Your turn}
\begin{itemize}
\raggedright
  \item \emph{Say it:} \emph{esconder}
  \item \emph{Say it:} \emph{esconder}, once more
  \item \emph{Say it:} a sentence of your own with \emph{esconder}
  \item \emph{Recall:} say the Spanish for the enigma, then the Spanish for to conceal, then say \emph{esconder} again
\end{itemize}

\begin{tcolorbox}[breakable,colback=teal!4,colframe=teal!35!black,title={Before you move on}]
What does esconder mean? (\textquotedblleft{}To hide\textquotedblright{}.) Say it once more.
\end{tcolorbox}
\section[el secreto]{el secreto — the secret; secret}

\label{lesson:ES-C532-secreto}



\begin{quote}
Before the new one: say the Spanish for to conceal, then the Spanish for to hide.
\end{quote}

\subsection*{You'll want to know: el secreto}
\textbf{el secreto} — \textquotedblleft{}the secret; secret\textquotedblright{}.

\begin{quote}
Es un secreto, no se lo digas a nadie.
\end{quote}

\begin{quote}
\emph{It is a secret, don't tell anyone.}
\end{quote}

As an adjective: \emph{una puerta secreta}.

\begin{cousinweb}
\emph{secreto} is a learned borrowing from Latin \emph{sēcrētus} 'set apart' (EtymDB), the participle of \emph{sēcernere} 'to separate', built on \emph{cernere} 'to sift'.
\end{cousinweb}

\begin{grammarlens}[title={What you've built}]
One more word for telling, explaining and describing at B1.
\end{grammarlens}

\subsection*{Your turn}
\begin{itemize}
\raggedright
  \item \emph{Say it:} \emph{el secreto}
  \item \emph{Say it:} \emph{el secreto}, once more
  \item \emph{Say it:} a sentence of your own with \emph{el secreto}
  \item \emph{Recall:} say the Spanish for to conceal, then the Spanish for to hide, then say \emph{el secreto} again
\end{itemize}

\begin{tcolorbox}[breakable,colback=teal!4,colframe=teal!35!black,title={Before you move on}]
What does el secreto mean? (\textquotedblleft{}The secret; secret\textquotedblright{}.) Say it once more.
\end{tcolorbox}
